\documentclass[10pt, a4paper]{article}
\usepackage[margin=0.5in]{geometry}
\usepackage{amsmath, amssymb, amsthm}
\usepackage{multicol}
\usepackage{enumitem}
\usepackage{titlesec}
\usepackage{xcolor}

% Compact formatting
\setlength{\parindent}{0pt}
\setlength{\parskip}{3pt}
\pagestyle{empty}

% Section styling
\titleformat{\section}{\large\bfseries\uppercase}{}{0em}{}[\hrule]
\titleformat{\subsection}{\bfseries\color{blue}}{}{0em}{}

% Environment definitions
\newtheorem*{theorem}{Theorem}
\newtheorem*{lemma}{Lemma}
\newtheorem*{definition}{Definition}

\begin{document}

\begin{center}
    {\Huge \textbf{Minimum Spanning Tree (MST) Cheat Sheet}}
\end{center}

\begin{multicols*}{2}

\section{Problem Definition}
Given a connected, undirected, weighted graph $G = (V, E)$ with weight function $w: E \rightarrow \mathbb{R}$.
\begin{itemize}[leftmargin=*]
    \item \textbf{Goal:} Find a spanning tree $T \subseteq E$ connecting all $|V|$ vertices such that the total weight $w(T) = \sum_{(u,v) \in T} w(u,v)$ is minimized.
    \item \textbf{Note:} If edge weights are distinct, the MST is unique.
\end{itemize}

\section{Generic MST Method}
An iterative approach to finding an MST.
\begin{itemize}[leftmargin=*]
    \item Maintain a subset of edges $A \subseteq E$.
    \item \textbf{Invariant:} $A$ is a \textbf{safe} subset (i.e., $A$ is a subset of some MST).
    \item \textbf{Step:} Find an edge $(u, v) \in E - A$ such that $A \cup \{(u, v)\}$ is still safe.
    \item \textbf{Loop:} Repeat until $|A| = |V| - 1$.
\end{itemize}

\section{Key Theorems \& Proofs}

\subsection{Definitions}
\begin{itemize}[leftmargin=*]
    \item \textbf{Cut $(S, V-S)$:} A partition of vertices $V$ into two disjoint sets.
    \item \textbf{Crosses:} An edge $(u, v)$ crosses the cut if one endpoint is in $S$ and the other in $V-S$.
    \item \textbf{Respects:} A cut respects set $A$ if no edge in $A$ crosses the cut.
    \item \textbf{Light Edge:} An edge crossing a cut with the minimum weight among all edges crossing that cut.
\end{itemize}

\subsection{The Cut Theorem}
\begin{theorem}
Let $A$ be a safe subset of $E$. Let $(S, V-S)$ be any cut that respects $A$. Let $(u, v)$ be a light edge crossing $(S, V-S)$. Then, $A \cup \{(u, v)\}$ is safe.
\end{theorem}

\textbf{Proof (Exchange Argument):}
\begin{enumerate}[leftmargin=*]
    \item Since $A$ is safe, there exists an MST $T$ such that $A \subseteq T$.
    \item If $(u, v) \in T$, we are done.
    \item If $(u, v) \notin T$, adding $(u, v)$ to $T$ creates a unique cycle.
    \item Since $(u, v)$ crosses the cut $(S, V-S)$, there must be at least one other edge $(x, y)$ on this cycle that also crosses the cut.
    \item Because $(u, v)$ is a light edge, $w(u, v) \le w(x, y)$.
    \item Construct $T' = (T \setminus \{(x, y)\}) \cup \{(u, v)\}$. $T'$ is a spanning tree.
    \item $w(T') = w(T) - w(x, y) + w(u, v) \le w(T)$.
    \item Since $T$ is an MST, $T'$ must also be an MST. Since $A \subseteq T$ and $(x, y) \notin A$ (because the cut respects $A$), $A \cup \{(u, v)\} \subseteq T'$. Thus, $A \cup \{(u, v)\}$ is safe. \qed
\end{enumerate}

\section{Kruskal's Algorithm}
Builds a forest that merges into a single tree. It is a greedy algorithm.

\subsection{Algorithm}
\begin{enumerate}[leftmargin=*]
    \item Sort all edges $E$ in non-decreasing order of weight.
    \item Initialize $A = \emptyset$. Create $|V|$ disjoint sets (one for each vertex).
    \item Iterate through sorted edges $(u, v)$:
    \item If $u$ and $v$ are in different sets (components):
    \begin{itemize}
        \item Add $(u, v)$ to $A$.
        \item Union the sets containing $u$ and $v$.
    \end{itemize}
\end{enumerate}

\subsection{Correctness}
Follows from the Cut Theorem. If $(u, v)$ connects two different trees $C_1$ and $C_2$, consider the cut $(V_{C_1}, V - V_{C_1})$. $(u, v)$ is the lightest edge crossing this cut (due to sorting), so it is safe to add.

\subsection{Data Structure: Disjoint Set Union (DSU)}
Operations: \texttt{Make-Set}, \texttt{Find-Set}, \texttt{Union}.
Using Path Compression and Union by Rank/Size.

\subsection{Time Complexity}
\begin{itemize}[leftmargin=*]
    \item \textbf{Sorting:} $\Theta(m \lg m)$ or $\Theta(m \lg n)$ (since $m < n^2$).
    \item \textbf{DSU Operations:} $O(m \alpha(n))$, where $\alpha$ is the inverse Ackermann function (nearly constant).
    \item \textbf{Total:} $\Theta(m \lg m)$ or $\Theta(m \lg n)$.
    \item If edges are pre-sorted: $O(m \alpha(n))$ (almost linear).
\end{itemize}
\textit{Notation: $n = |V|, m = |E|$.}

\section{Prim's Algorithm}
Grows a single tree from a root vertex $r$. Also greedy.

\subsection{Algorithm}
\begin{enumerate}[leftmargin=*]
    \item Initialize $S = \{r\}$ (vertices in MST), $A = \emptyset$.
    \item Maintain a priority queue $Q$ of vertices in $V-S$, keyed by the minimum weight edge connecting them to $S$.
    \item $key[v] = \min \{w(u, v) \mid u \in S\}$. If no edge, $\infty$.
    \item While $Q$ is not empty:
    \begin{itemize}
        \item Extract min vertex $u$ from $Q$.
        \item Add $u$ to $S$, add edge $(u.\pi, u)$ to $A$.
        \item Update keys of neighbors $v$ of $u$: if $w(u, v) < key[v]$, update $key[v] = w(u, v)$ and set parent $\pi[v] = u$.
    \end{itemize}
\end{enumerate}

\subsection{Correctness}
Maintains invariant that $A$ is a safe set forming a tree on $S$. At each step, it selects the light edge crossing cut $(S, V-S)$, which is safe by the Cut Theorem.

\subsection{Time Complexity}
Depends on Priority Queue (PQ) implementation:
\begin{center}
\begin{tabular}{|l|l|l|}
\hline
\textbf{PQ Type} & \textbf{Operations} & \textbf{Total Time} \\ \hline
Binary Heap & $O(n)$ Extract-Min & $O(m \lg n)$ \\
 & $O(m)$ Decrease-Key & \\ \hline
Fibonacci Heap & $O(n)$ Extract-Min & $O(m + n \lg n)$ \\
 & $O(1)$ amortized Dec-Key & \\ \hline
Adjacency Matrix & Array search & $O(n^2)$ \\ \hline
\end{tabular}
\end{center}
\textit{Note: Binary Heap is standard. Fibonacci is theoretical optimum for dense graphs.}

\section{Comparison Summary}
\begin{center}
\renewcommand{\arraystretch}{1.2}
\begin{tabular}{|p{0.25\linewidth}|p{0.3\linewidth}|p{0.3\linewidth}|}
\hline
\textbf{Feature} & \textbf{Kruskal's} & \textbf{Prim's} \\ \hline
\textbf{Approach} & Process edges (global) & Process vertices (local/grow) \\ \hline
\textbf{Structure} & Maintains Forest & Maintains Single Tree \\ \hline
\textbf{Best For} & Sparse Graphs & Dense Graphs (w/ Fib Heap) \\ \hline
\textbf{Complexity} & $O(m \lg n)$ & $O(m \lg n)$ or $O(m+n \lg n)$ \\ \hline
\end{tabular}
\end{center}

\end{multicols*}
\end{document}